% CVPR 2027 Paper Template; see https://github.com/cvpr-org/author-kit

\documentclass[10pt,twocolumn,letterpaper]{article}

%%%%%%%%% PAPER TYPE  - PLEASE UPDATE FOR FINAL VERSION
% \usepackage{cvpr}              % To produce the CAMERA-READY version
\usepackage[review]{cvpr}      % To produce the REVIEW version
% \usepackage[pagenumbers]{cvpr} % To force page numbers, e.g. for an arXiv version

% Import additional packages in the preamble file, before hyperref
%% This file contains a number of tweaks that are typically applied to the main document.
%% They are not enabled by default, but can be enabled by uncommenting the relevant lines.

%%
%% Inline annotations; for predefined colors, refer to "dvipsnames" in the xcolor package:
%% https://tinyurl.com/overleaf-colors
%%
\newcommand{\red}[1]{{\color{red}#1}}
\newcommand{\todo}[1]{{\color{red}#1}}
\newcommand{\TODO}[1]{\textbf{\color{red}[TODO: #1]}}
%%
%% disable for camera ready / submission by uncommenting these lines  
%%
% \renewcommand{\TODO}[1]{}
% \renewcommand{\todo}[1]{#1}

%%
%% work harder in optimizing text layout. Typically shrinks text by 1/6 of page, enable
%% it at the very end of the writing process, when you are just above the page limit
%%
% \usepackage{microtype}

%%
%% fine-tune paragraph spacing
%%
% \renewcommand{\paragraph}[1]{\vspace{.5em}\noindent\textbf{#1.}}

%%
%% globally adjusts space between figure and caption
%%
% \setlength{\abovecaptionskip}{.5em}


%%
%% Allows "the use of \paper to refer to the project name"
%% with automatic management of space at the end of the word
%%
% \usepackage{xspace}
% \newcommand{\paper}{ProjectName\xspace}

%%
%% Commonly used math definitions
%%
% \DeclareMathOperator*{\argmin}{arg\,min}
% \DeclareMathOperator*{\argmax}{arg\,max}

%%
%% Tigthen underline
%%
% \usepackage{soul}
% \setuldepth{foobar}

% It is strongly recommended to use hyperref, especially for the review version.
% hyperref with option pagebackref eases the reviewers' job.
% Please disable hyperref *only* if you encounter grave issues, 
% e.g. with the file validation for the camera-ready version.
%
% If you comment hyperref and then uncomment it, you should delete *.aux before re-running LaTeX.
% (Or just hit 'q' on the first LaTeX run, let it finish, and you should be clear).
\definecolor{cvprblue}{rgb}{0.21,0.49,0.74}
\usepackage[pagebackref,breaklinks,colorlinks,allcolors=cvprblue]{hyperref}

\def\confName{CVPR}
\def\confYear{2027}

%%%%%%%%% TITLE - PLEASE UPDATE
\title{CVPR 2027 - Vision Anomaly Detection}

%%%%%%%%% AUTHORS - PLEASE UPDATE
\author{First Author\\
Institution1\\
Institution1 address\\
{\tt\small firstauthor@i1.org}
% For a paper whose authors are all at the same institution,
% omit the following lines up until the closing ``}''.
% Additional authors and addresses can be added with ``\and'',
% just like the second author.
% To save space, use either the email address or home page, not both
\and
Second Author\\
Institution2\\
First line of institution2 address\\
{\tt\small secondauthor@i2.org}
}

\begin{document}
\maketitle

% All paper sections are written directly in this file.
% Replace the TODOs with paper content as the research develops.
\begin{abstract}
\TODO{Summarize the problem, the proposed approach, the main experimental
findings, and their significance.}
\end{abstract}

%-------------------------------------------------------------------------
\section{Introduction}
\label{sec:intro}

\TODO{Introduce the problem and its motivation, explain the gap in existing
approaches, and summarize the proposed solution and main contributions.}

%-------------------------------------------------------------------------
\section{Related Work}
\label{sec:related_work}

\TODO{Review the most relevant research, organize it by topic, and explain
how the proposed approach relates to prior work.}

%-------------------------------------------------------------------------
\section{Method}
\label{sec:method}

\subsection{Problem Formulation}
\label{sec:problem_formulation}

\TODO{Define the task, notation, inputs, outputs, and assumptions.}

\subsection{Approach}
\label{sec:approach}

\TODO{Describe the method and its main components, including the training
objective and inference procedure where applicable.}

%-------------------------------------------------------------------------
\section{Experiments}
\label{sec:experiments}

\subsection{Experimental Setup}
\label{sec:experimental_setup}

\TODO{Describe datasets, evaluation protocols, metrics, baselines, and
implementation details needed to reproduce the experiments.}

\subsection{Main Results}
\label{sec:main_results}

\TODO{Report measured results, compare with baselines, and discuss the
findings. Add numerical claims only after the experiments are completed.}

\subsection{Ablation Studies}
\label{sec:ablations}

\TODO{Evaluate the contribution of individual components and sensitivity
to key design choices.}

\subsection{Qualitative Results and Limitations}
\label{sec:qualitative_results}

\TODO{Present representative examples, analyze failure cases, and discuss
the limitations of the approach.}

%-------------------------------------------------------------------------
\section{Conclusion}
\label{sec:conclusion}

\TODO{Summarize the contributions and supported findings, and identify
directions for future work.}

%-------------------------------------------------------------------------
% Add real references to main.bib and cite them in the sections above.
{
    \small
    \bibliographystyle{ieeenat_fullname}
    \bibliography{main}
}

% Optional supplementary material can be written here, in this same file.
% Exclude it from the main-paper submission when required.
% \clearpage
% \setcounter{page}{1}
% \maketitlesupplementary
% \appendix
% \section{Additional Implementation Details}
% \label{sec:supplementary}

\end{document}
